% !TEX program = lualatex
\documentclass[aspectratio=169,professionalfonts]{beamer}
\newcommand{\slidemode}{1}  % カラー投影モード
\usepackage{beamer_template}
\usepackage{enumerate}

\newcommand{\LectureNum}{1}
\newcommand{\LectureTitle}{ラプラス変換の定義}
\newcommand{\CourseName}{応用数学}
\renewcommand{\TermName}{前期}

\begin{document}

\begin{frame}{第1回　ラプラス変換の定義}
  \begin{block}{本回の内容}
  ラプラス変換の定義
  \end{block}
  \vfill\hfill{\scriptsize \textcolor{gray}{第2章 ラプラス変換　§1}}
\end{frame}

\begin{frame}{定義：ラプラス変換}
  \begin{block}{}
  \[
    F(s) = \int_{0}^\infty e^{-st} f(t) dt = \lim_{T \to \infty} \int_{0}^T e^{-st} f(t) dt
  \]
  \end{block}
  \vfill\hfill{\scriptsize \textcolor{gray}{p.\,46--49}}
\end{frame}

\begin{frame}{例題1}
  \begin{exampleblock}{問題}
    $f(t) = 1$ のラプラス変換を求めよ\\[2pt]
  \end{exampleblock}
\end{frame}

\begin{frame}{例題1　解答}
  \begin{exampleblock}{解答}
  \begin{enumerate}[1.]
    \item[] $\lim_{t \to \infty} e^{-st} = 0$ だから
    \item[] $F(s) = \int_{0}^\infty e^{-st} dt = [-1/s \cdot e^{-st}]_{0}^\infty = 1/s$
    \item[] $\int_{0}^\infty dt$ は存在しない
  \end{enumerate}
  \end{exampleblock}
  \begin{alertblock}{答}
    $L[1] = 1/s  (s > 0)$
  \end{alertblock}
\end{frame}

\begin{frame}{例題2}
  \begin{exampleblock}{問題}
    $f(t) = t$ のラプラス変換を求めよ\\[2pt]
  \end{exampleblock}
\end{frame}

\begin{frame}{例題2　解答}
  \begin{exampleblock}{解答}
  \begin{enumerate}[1.]
    \item[] 分部積分を用いると
    \item[] $F(s) = \int_{0}^\infty e^{-st}\cdot t dt = [-1/s \cdot e^{-st} \cdot t]_{0}^\infty + 1/s \int_{0}^\infty e^{-st} dt$
    \item[] ロピタルの定理より $\lim_{t \to \infty} t\cdot e^{-st} = \lim_{t \to \infty} t/e^{st} = \lim_{t \to \infty} 1/(s\cdot e^{st}) = 0$
    \item[] よって $F(s) = 0 + (1/s)\cdot L[1] = (1/s)\cdot(1/s) = 1/s^{2}$
    \item[] $\int_{0}^\infty t dt$ は存在しない
  \end{enumerate}
  \end{exampleblock}
  \begin{alertblock}{答}
    $L[t] = 1/s^{2}  (s > 0)$
  \end{alertblock}
\end{frame}

\begin{frame}{例題3}
  \begin{exampleblock}{問題}
    $f(t) = e^{at}$ （ a は定数）のラプラス変換を求めよ\\[2pt]
  \end{exampleblock}
\end{frame}

\begin{frame}{例題3　解答}
  \begin{exampleblock}{解答}
  \begin{enumerate}[1.]
    \item[] $F(s) = \int_{0}^\infty e^{-st} \cdot e^{at} dt = \int_{0}^\infty e^{-(s-a)t} dt$
    \item[] 例題1 と同様に、s-a > 0 すなわち s > a のとき
    \item[] $F(s) = [-1/(s-a) \cdot e^{-(s-a)t}]_{0}^\infty = 1/(s-a)$
    \item[] $e^{-(s-a)t}$ は存在しない
  \end{enumerate}
  \end{exampleblock}
  \begin{alertblock}{答}
    $L[e^{at}] = 1/(s-a)  (s > a)$
  \end{alertblock}
\end{frame}

\begin{frame}{例題4}
  \begin{exampleblock}{問題}
    $f(t) = \sin t$ のラプラス変換を求めよ\\[2pt]
  \end{exampleblock}
\end{frame}

\begin{frame}{例題4　解答}
  \begin{exampleblock}{解答}
  \begin{enumerate}[1.]
    \item[] s > 0 のとき $\lim_{t \to \infty} e^{-st}\sin t = 0, \lim_{t \to \infty} e^{-st}\cos t = 0$
    \item[] 部分積分を2回繰り返すと
    \item[] $F(s) = [-e^{-st}(-\cos t)]_{0}^\infty - s\int_{0}^\infty e^{-st}(-\cos t)dt$
    \item[] $= 1 - s\cdot L[\cos t]$
    \item[] 同様に $L[\cos t] = s\cdot F(s)$ から代入すると
    \item[] $F(s) = 1 - s^{2}F(s)   \to   (1 + s^{2})F(s) = 1$
  \end{enumerate}
  \end{exampleblock}
  \begin{alertblock}{答}
    $L[\sin t] = 1/(s^{2}+1)  (s > 0)$
  \end{alertblock}
\end{frame}

\begin{frame}{例題6}
  \begin{exampleblock}{問題}
    正の実数 a, b （ただし a < b ）について、関数 $f(t)$ を\\[2pt]
    $f(t) = \begin{cases} 0 & (0 \leq t < a) \\ 1 & (a \leq t < b) \\ 0 & (b \leq t) \end{cases}$\\[2pt]
    で定めるとき、 $f(t)$ を単位ステップ関数を用いて表し、 $L[f(t)]$ を求めよ\\[2pt]
  \end{exampleblock}
\end{frame}

\begin{frame}{例題6　解答}
  \begin{exampleblock}{解答}
  \begin{enumerate}[1.]
    \item[] 単位ステップ関数 U(t-a), U(t-b) を用いると
    \item[] $f(t) = U(t-a) - U(t-b)$
    \item[] 線形性より
    \item[] $L[f(t)] = L[U(t-a)] - L[U(t-b)]$
    \item[] $= e^{-as}/s - e^{-bs}/s$
  \end{enumerate}
  \end{exampleblock}
  \begin{alertblock}{答}
    $L[f(t)] = (e^{-as} - e^{-bs})/s  (s > 0)$
  \end{alertblock}
\end{frame}

\begin{frame}{演習問題（p.\,46--49）}
  \begin{exampleblock}{自分で解いてみよう}
  \begin{description}
    \item[問1] $L[3 + 2^t]$ を線形性を用いて求めよ
    \item[問2] $L[3 + 2^t]$ を求めよ
    \item[問3] $L[e^{-t} + e^{-2t}]$ を求めよ
    \item[問4] 次の公式を証明せよ $: L[\cos t] = s/(s^{2}+1) (s > 0)$
    \item[問5] 双曲線関数 sinh t, cosh t のラプラス変換を求めよ
    \item[問6] 関数 $g(t) = U(t-2)$ のグラフをかけ。また $L[U(t-2)]$ を求めよ
    \item[問7] 正の実数 a, b, c （ただし a < b < c ）について，関数 $f(t)$ を
    $f(t) = { 1  (0 < t < a, b \leq t < c),  0  (a \leq t < b, t \geq c) }$\\
    で定めるとき， $f(t)$ を単位ステップ関数を用いて表せ。また， $L[f(t)]$ を求めよ。\\
  \end{description}
  \end{exampleblock}
\end{frame}

\end{document}
